% Gerado por scripts/migrate_markdown_to_latex.py.
% Fonte migrada: fases/01_validacao_conceitual/docs/validacao_transformer.md


\chapter{Validacao arquitetural da NN usada no estudo}




Este documento registra se a rede neural usada no estudo pode ser chamada de Transformer.



Ele tambem registra a ponte metodologica entre:



\begin{Verbatim}[fontsize=\small,breaklines=true,label=text]
Vaswani et al. (2017) -> propriedade esperada -> script Python -> teste deterministico -> limite da conclusao
\end{Verbatim}






Essa ponte e importante porque o estudo nao deve depender do nome de uma classe ou de uma afirmacao informal. Para dizer que um trecho e de Transformer, precisamos mostrar que ele possui a propriedade descrita no artigo e que essa propriedade foi testada.



\section{Resultado atual}




A implementacao \texttt{validacao.\allowbreak{}transformer.\allowbreak{}Simplified\allowbreak{}Transformer\allowbreak{}Block} esta validada como:



\begin{Verbatim}[fontsize=\small,breaklines=true,label=text]
bloco_transformer_simplificado
\end{Verbatim}




Ela pertence a familia Transformer porque passa nos criterios essenciais de um bloco:


\begin{itemize}
\item recebe entrada sequencial \texttt{[batch, seq\_\allowbreak{}len, d\_\allowbreak{}model]};
\item preserva shape sequencial na saida;
\item projeta a entrada em Q, K e V;
\item usa atencao escalonada compativel com \texttt{softmax(QK\textasciicircum{}T /\allowbreak{} sqrt(d\_\allowbreak{}k))V};
\item usa multi-head attention;
\item mistura informacao entre posicoes;
\item usa informacao posicional sinusoidal;
\item possui residual, normalizacao e feed-forward;
\item e deterministica em modo de avaliacao.
\end{itemize}




Ela nao deve ser chamada de 'Transformer completo' no artigo, porque o recorte experimental usa um bloco controlado, nao uma arquitetura completa com pilha de camadas, cabeca de tarefa, treinamento e protocolo de modelo final.



\section{Criterio de rastreabilidade}



Cada componente aceito no estudo precisa passar por cinco perguntas:


\begin{enumerate}
\item O componente existe no Transformer de Vaswani et al.?
\item Qual propriedade observavel esse componente deve ter?
\item Onde essa propriedade aparece no nosso codigo?
\item Qual teste deterministico verifica essa propriedade?
\item O que esse teste ainda nao permite concluir?
\end{enumerate}



Assim, a validacao nao diz apenas 'o codigo roda'. Ela diz 'esta operacao corresponde a este trecho definido pela literatura, dentro destes limites'.



\section{Matriz Vaswani -> codigo -> teste}



\begin{landscape}
\scriptsize
\begin{longtable}{@{}>{\RaggedRight\arraybackslash}p{0.117\linewidth}>{\RaggedRight\arraybackslash}p{0.117\linewidth}>{\RaggedRight\arraybackslash}p{0.117\linewidth}>{\RaggedRight\arraybackslash}p{0.117\linewidth}>{\RaggedRight\arraybackslash}p{0.117\linewidth}>{\RaggedRight\arraybackslash}p{0.117\linewidth}>{\RaggedRight\arraybackslash}p{0.117\linewidth}@{}}
\toprule
Trecho no Transformer de Vaswani & Propriedade esperada & Onde esta no codigo & Teste deterministico & Comparacao usada & O que o teste prova & Limite metodologico \\
\midrule
\endfirsthead
\toprule
Trecho no Transformer de Vaswani & Propriedade esperada & Onde esta no codigo & Teste deterministico & Comparacao usada & O que o teste prova & Limite metodologico \\
\midrule
\endhead
Entrada sequencial & O modulo recebe tensores organizados como sequencia, com dimensao de lote, posicao e representacao. & \texttt{Simplified\allowbreak{}Transformer\allowbreak{}Block.\allowbreak{}forward}, em \texttt{fases/\allowbreak{}01\_\allowbreak{}validacao\_\allowbreak{}conceitual/\allowbreak{}validacao/\allowbreak{}transformer.\allowbreak{}py}. & \texttt{test\_\allowbreak{}simplified\_\allowbreak{}transformer\_\allowbreak{}block\_\allowbreak{}has\_\allowbreak{}expected\_\allowbreak{}components\_\allowbreak{}and\_\allowbreak{}shape}. & Shape esperado \texttt{[batch, seq\_\allowbreak{}len, d\_\allowbreak{}model]}. & O bloco recebe e devolve uma representacao sequencial. & Nao prova tokenizacao, embedding real ou tarefa de linguagem. \\
Projecoes lineares densas & Entradas sao transformadas por matrizes aprendidas; no attention original essas projecoes geram Q, K, V e saida. & \texttt{q\_\allowbreak{}proj}, \texttt{k\_\allowbreak{}proj}, \texttt{v\_\allowbreak{}proj}, \texttt{out\_\allowbreak{}proj} em \texttt{fases/\allowbreak{}01\_\allowbreak{}validacao\_\allowbreak{}conceitual/\allowbreak{}validacao/\allowbreak{}transformer.\allowbreak{}py}; \texttt{dense\_\allowbreak{}projection\_\allowbreak{}torch} em \texttt{fases/\allowbreak{}01\_\allowbreak{}validacao\_\allowbreak{}conceitual/\allowbreak{}validacao/\allowbreak{}dense.\allowbreak{}py}. & \texttt{test\_\allowbreak{}dense\_\allowbreak{}projection\_\allowbreak{}matches\_\allowbreak{}manual\_\allowbreak{}definition}; auditoria de Q/K/V. & Formula \texttt{Y = XW\textasciicircum{}T + b} calculada manualmente. & A operacao linear usada no bloco segue a definicao matematica esperada. & Nao prova que os pesos treinados sao bons; os pesos do bloco controlado sao fixos/iniciais. \\
Projecoes Q, K e V & A mesma sequencia gera consultas, chaves e valores para self-attention. & \texttt{project\_\allowbreak{}qkv} em \texttt{fases/\allowbreak{}01\_\allowbreak{}validacao\_\allowbreak{}conceitual/\allowbreak{}validacao/\allowbreak{}transformer.\allowbreak{}py}. & \texttt{test\_\allowbreak{}transformer\_\allowbreak{}audit\_\allowbreak{}validates\_\allowbreak{}current\_\allowbreak{}nn\_\allowbreak{}as\_\allowbreak{}simplified\_\allowbreak{}transformer\_\allowbreak{}block}. & Inspecao estrutural: existencia de \texttt{q\_\allowbreak{}proj}, \texttt{k\_\allowbreak{}proj}, \texttt{v\_\allowbreak{}proj}. & O bloco tem as tres projecoes exigidas para attention de Transformer. & Nao prova arquitetura completa encoder-decoder. \\
Scaled dot-product attention & O nucleo de atencao deve calcular \texttt{softmax(QK\textasciicircum{}T /\allowbreak{} sqrt(d\_\allowbreak{}k))V}. & \texttt{scaled\_\allowbreak{}dot\_\allowbreak{}product\_\allowbreak{}attention\_\allowbreak{}torch} e \texttt{scaled\_\allowbreak{}dot\_\allowbreak{}product\_\allowbreak{}attention\_\allowbreak{}numpy} em \texttt{fases/\allowbreak{}01\_\allowbreak{}validacao\_\allowbreak{}conceitual/\allowbreak{}validacao/\allowbreak{}attention.\allowbreak{}py}; adaptador em \texttt{validacao/\allowbreak{}comparacao\_\allowbreak{}frameworks.\allowbreak{}py}. & Testes de referencia NumPy/PyTorch e \texttt{test\_\allowbreak{}all\_\allowbreak{}frameworks\_\allowbreak{}match\_\allowbreak{}in\_\allowbreak{}every\_\allowbreak{}case}. & Valor manual pequeno, NumPy independente, \texttt{torch.\allowbreak{}nn.\allowbreak{}functional.\allowbreak{}scaled\_\allowbreak{}dot\_\allowbreak{}product\_\allowbreak{}attention} e \texttt{keras.\allowbreak{}ops.\allowbreak{}nn.\allowbreak{}dot\_\allowbreak{}product\_\allowbreak{}attention} com backend TensorFlow. & O algoritmo bate com a formula e com APIs publicas de dois frameworks nos mesmos tensores. & Nao prova por si so que existe um bloco completo nem compara desempenho; prova o nucleo de attention. \\
Mascara de atencao & A atencao pode receber restricoes aditivas, como ocorre em variantes mascaradas. & Parametro \texttt{additive\_\allowbreak{}attention\_\allowbreak{}mask} em \texttt{fases/\allowbreak{}01\_\allowbreak{}validacao\_\allowbreak{}conceitual/\allowbreak{}validacao/\allowbreak{}attention.\allowbreak{}py}. & \texttt{test\_\allowbreak{}scaled\_\allowbreak{}dot\_\allowbreak{}product\_\allowbreak{}attention\_\allowbreak{}additive\_\allowbreak{}mask\_\allowbreak{}matches\_\allowbreak{}pytorch\_\allowbreak{}reference}. & \texttt{torch.\allowbreak{}nn.\allowbreak{}functional.\allowbreak{}scaled\_\allowbreak{}dot\_\allowbreak{}product\_\allowbreak{}attention} com \texttt{attn\_\allowbreak{}mask}. & A mascara usada no nosso nucleo e compativel com a referencia do PyTorch. & O bloco experimental padrao nao usa mascara causal por padrao. \\
Dependencia entre posicoes & A saida de uma posicao pode depender de outra posicao da sequencia. & Self-attention em \texttt{fases/\allowbreak{}01\_\allowbreak{}validacao\_\allowbreak{}conceitual/\allowbreak{}validacao/\allowbreak{}attention.\allowbreak{}py} e auditoria em \texttt{fases/\allowbreak{}01\_\allowbreak{}validacao\_\allowbreak{}conceitual/\allowbreak{}validacao/\allowbreak{}auditoria\_\allowbreak{}transformer.\allowbreak{}py}. & \texttt{test\_\allowbreak{}attention\_\allowbreak{}output\_\allowbreak{}depends\_\allowbreak{}on\_\allowbreak{}other\_\allowbreak{}sequence\_\allowbreak{}positions}; \texttt{test\_\allowbreak{}transformer\_\allowbreak{}audit\_\allowbreak{}validates\_\allowbreak{}current\_\allowbreak{}nn\_\allowbreak{}as\_\allowbreak{}simplified\_\allowbreak{}transformer\_\allowbreak{}block}. & Alteracao controlada de um token e observacao de mudanca em outra posicao. & A operacao nao trata cada posicao de forma isolada; ha mistura de informacao sequencial. & Nao mede qualidade semantica dessa dependencia em uma tarefa real. \\
Multi-head attention & A representacao e dividida em cabecas, cada cabeca aplica attention, e depois as cabecas sao recombinadas. & \texttt{split\_\allowbreak{}heads}, \texttt{combine\_\allowbreak{}heads} e \texttt{multi\_\allowbreak{}head\_\allowbreak{}attention\_\allowbreak{}from\_\allowbreak{}qkv} em \texttt{fases/\allowbreak{}01\_\allowbreak{}validacao\_\allowbreak{}conceitual/\allowbreak{}validacao/\allowbreak{}attention.\allowbreak{}py}. & \texttt{test\_\allowbreak{}multi\_\allowbreak{}head\_\allowbreak{}attention\_\allowbreak{}preserves\_\allowbreak{}shape\_\allowbreak{}and\_\allowbreak{}matches\_\allowbreak{}split\_\allowbreak{}reference}. & Aplicacao da SDPA do PyTorch em cada head e recombinacao manual. & A mecanica de separar/aplicar/recombinar heads esta correta. & Nao prova que o numero de heads escolhido e otimo. \\
Informacao posicional & Como attention pura nao conhece ordem por si so, o bloco precisa receber algum sinal de posicao. & \texttt{sinusoidal\_\allowbreak{}position\_\allowbreak{}encoding} e \texttt{add\_\allowbreak{}positional\_\allowbreak{}information} em \texttt{fases/\allowbreak{}01\_\allowbreak{}validacao\_\allowbreak{}conceitual/\allowbreak{}validacao/\allowbreak{}transformer.\allowbreak{}py}. & Auditoria: \texttt{report.\allowbreak{}has\_\allowbreak{}positional\_\allowbreak{}information\_\allowbreak{}or\_\allowbreak{}justification}. & Criterio arquitetural derivado de Vaswani et al. & O bloco adiciona informacao posicional antes das projecoes. & Nao prova que essa codificacao e a melhor para todas as tarefas. \\
Residual, normalizacao e FFN & O bloco Transformer usa subcamadas com conexoes residuais, normalizacao e rede feed-forward. & \texttt{norm1}, \texttt{norm2}, \texttt{ffn} e \texttt{out\_\allowbreak{}proj} em \texttt{fases/\allowbreak{}01\_\allowbreak{}validacao\_\allowbreak{}conceitual/\allowbreak{}validacao/\allowbreak{}transformer.\allowbreak{}py}. & \texttt{test\_\allowbreak{}simplified\_\allowbreak{}transformer\_\allowbreak{}block\_\allowbreak{}has\_\allowbreak{}expected\_\allowbreak{}components\_\allowbreak{}and\_\allowbreak{}shape}; auditoria estrutural. & Checklist arquitetural do bloco. & O bloco tem os componentes essenciais para ser classificado como bloco Transformer simplificado. & Ainda falta pilha completa de camadas e cabeca de tarefa para chamar de Transformer completo. \\
Inferencia deterministica & Em avaliacao, a mesma entrada deve produzir a mesma saida, evitando ruido experimental. & \texttt{audit\_\allowbreak{}transformer\_\allowbreak{}module} em \texttt{fases/\allowbreak{}01\_\allowbreak{}validacao\_\allowbreak{}conceitual/\allowbreak{}validacao/\allowbreak{}auditoria\_\allowbreak{}transformer.\allowbreak{}py}. & \texttt{test\_\allowbreak{}transformer\_\allowbreak{}audit\_\allowbreak{}validates\_\allowbreak{}current\_\allowbreak{}nn\_\allowbreak{}as\_\allowbreak{}simplified\_\allowbreak{}transformer\_\allowbreak{}block}. & Duas execucoes em \texttt{eval()} com a mesma entrada. & O nucleo pode ser usado em benchmarks controlados. & Nao substitui avaliacao estatistica de desempenho em varias repeticoes. \\
Controle negativo & Uma NN comum nao deve ser aceita como Transformer so por ser rede neural. & Auditoria aplicada a \texttt{torch.\allowbreak{}nn.\allowbreak{}Sequential(Linear, ReLU, Linear)}. & \texttt{test\_\allowbreak{}transformer\_\allowbreak{}audit\_\allowbreak{}rejects\_\allowbreak{}plain\_\allowbreak{}neural\_\allowbreak{}network\_\allowbreak{}as\_\allowbreak{}transformer}. & Ausencia de Q/K/V e de scaled attention. & O validador distingue uma NN generica de um bloco da familia Transformer. & Nao classifica todos os tipos possiveis de arquitetura; classifica o recorte do estudo. \\
\bottomrule
\end{longtable}
\end{landscape}



\section{O que esta comprovado}





Os testes comprovam que os scripts Python implementam deterministicamente os trechos centrais escolhidos do Transformer de Vaswani et al. para o escopo do projeto:


\begin{itemize}
\item projecao linear densa;
\item geracao de Q, K e V;
\item scaled dot-product attention;
\item mascara aditiva de attention;
\item dependencia entre posicoes da sequencia;
\item multi-head attention;
\item informacao posicional sinusoidal;
\item subcamadas residual, normalizacao e FFN;
\item comportamento deterministico em inferencia;
\item rejeicao de uma NN comum como controle negativo.
\end{itemize}


Portanto, o objeto experimental atual e valido para estudar:



\begin{Verbatim}[fontsize=\small,breaklines=true,label=text]
operacoes centrais de Transformer
\end{Verbatim}



e tambem:



\begin{Verbatim}[fontsize=\small,breaklines=true,label=text]
bloco Transformer simplificado
\end{Verbatim}



\section{O que nao esta comprovado}





Os testes nao comprovam que temos a arquitetura completa do artigo \emph{Attention Is All You Need}. Para isso, seria necessario incluir e validar elementos fora do recorte atual, como:


\begin{itemize}
\item pilha completa de blocos;
\item organizacao completa de encoder e/ou decoder;
\item embeddings de entrada e saida;
\item cabeca de tarefa;
\item treinamento ou carregamento de pesos de um modelo real;
\item avaliacao em tarefa ou dataset.
\end{itemize}


Por conta disso, a frase metodologicamente correta e:



\begin{Verbatim}[fontsize=\small,breaklines=true,label=text]
O estudo avalia trechos computacionalmente centrais de um bloco Transformer
simplificado, derivados da arquitetura de Vaswani et al. (2017), com foco em
inferencia controlada.
\end{Verbatim}



E a frase que deve ser evitada neste momento e:



\begin{Verbatim}[fontsize=\small,breaklines=true,label=text]
O estudo avalia uma Transformer completa.
\end{Verbatim}



\section{Como a verificacao e feita}




A auditoria esta implementada em \texttt{validacao.\allowbreak{}auditoria\_\allowbreak{}transformer}. Ela executa testes comportamentais e estruturais:


\begin{enumerate}
\item passa uma sequencia pelo modulo;
\item verifica se a saida preserva \texttt{[batch, seq\_\allowbreak{}len, d\_\allowbreak{}model]};
\item inspeciona Q/K/V, output projection, norm e FFN;
\item compara a atencao interna com \texttt{torch.\allowbreak{}nn.\allowbreak{}functional.\allowbreak{}scaled\_\allowbreak{}dot\_\allowbreak{}product\_\allowbreak{}attention};
\item altera um token e verifica se outra posicao pode mudar;
\item roda a mesma entrada duas vezes em \texttt{eval()} para testar determinismo;
\item aplica a classificacao \texttt{nao\_\allowbreak{}aderente}, \texttt{operacao\_\allowbreak{}inspirada\_\allowbreak{}em\_\allowbreak{}transformer},
\texttt{bloco\_\allowbreak{}transformer\_\allowbreak{}simplificado} ou \texttt{transformer}.
\end{enumerate}


\section{Controle negativo}





A mesma auditoria rejeita uma rede neural comum composta por camadas lineares e ReLU. Esse controle e importante: ele mostra que o teste nao aceita qualquer NN como Transformer.



\section{Regra para o texto final}



No artigo, usar:


\begin{itemize}
\item 'bloco Transformer simplificado' para o nucleo experimental atual;
\item 'operacao inspirada em Transformer' quando o teste isolar attention,
projecao densa ou KV cache;
\item 'Transformer' apenas se houver arquitetura completa ou justificativa formal
para todos os componentes ausentes.
\end{itemize}
